\section{Experimental Setup}
\label{sec:experiments}

\subsection{Benchmarks and split provenance}

Our suite contains 12 datasets spanning clinical detection (CHB-MIT, TUAB,
TUEV), sleep staging (ISRUC, HMC), emotion recognition (FACED, SEED-V), motor
imagery (PhysioNet-MI, SHU-MI), imagined speech (BCIC2020-3), depression
detection (Mumtaz2016), and mental workload (MentalArithmetic). Table~
\ref{tab:datasets} gives the input shapes and split sizes. All recordings are
represented at 200\,Hz in the common preprocessing pipeline. Regular loaders
use the $[B,C,T]$ convention; ISRUC preserves its sequence of 20 labelled
epochs.

\begin{table*}[t]
\centering
\caption{The 12-dataset evaluation suite and inherited benchmark split sizes.}
\label{tab:datasets}
\small
\setlength{\tabcolsep}{3pt}
\renewcommand{\arraystretch}{0.96}
\begin{tabular*}{\textwidth}{@{\extracolsep{\fill}}lccrrr@{}}
\toprule
Dataset & Input shape & Classes & Train & Val & Test \\
\midrule
\familyrow{6}{Clinical detection}
CHB-MIT & $16\times2000$ & 2 & 286,964 & 23,065 & 16,964 \\
TUAB & $16\times2000$ & 2 & 297,103 & 75,407 & 36,945 \\
TUEV & $16\times1000$ & 6 & 68,445 & 15,487 & 29,421 \\
\familyrow{6}{Sleep staging}
ISRUC & $20\times6\times6000$ & 5 & 3,559 & 468 & 435 \\
HMC & $4\times6000$ & 5 & 91,248 & 22,124 & 23,871 \\
\familyrow{6}{Affective computing}
FACED & $32\times2000$ & 9 & 6,720 & 1,680 & 1,932 \\
SEED-V & $62\times200$ & 5 & 34,432 & 42,960 & 40,352 \\
\familyrow{6}{Brain--computer interfaces}
PhysioNet-MI & $64\times800$ & 4 & 6,300 & 1,734 & 1,803 \\
SHU-MI & $32\times800$ & 2 & 7,210 & 2,431 & 2,347 \\
BCIC2020-3 & $64\times600$ & 5 & 4,500 & 750 & 750 \\
\familyrow{6}{Clinical and cognitive state decoding}
Mumtaz2016 & $19\times1000$ & 2 & 4,891 & 1,041 & 1,211 \\
MentalArithmetic & $20\times1000$ & 2 & 1,343 & 172 & 192 \\
\bottomrule
\end{tabular*}
\tablenote{Counts are local input segments, except ISRUC (20-epoch sequences)
and HMC (30-second epochs). Shapes omit the batch dimension; classes are
semantic classes. TUEV counts use v2.0.1.}
\end{table*}

For all 12 datasets, we follow the CBraMod/REVE benchmark definitions for
record segmentation, sample construction, preprocessing order, and public
train/validation/test partitions \citep{wang2024cbramod,elouahidi2025reve}.
Our implementation preserves these conventions; the explicit study-specific
operation is the final clipping and scaling applied before folding (specified
below). For example, CHB-MIT uses
CHB01--20/CHB21--22/CHB23--24 for train/validation/test. One version caveat is
TUEV: our preprocessing uses v2.0.1 rather than the v2.0.0 release used by one
reference implementation. Both releases share the official 29,421-sample test
set, so the test partition is unchanged; only the training and validation
counts differ. TUEV comparisons should therefore be interpreted with this
release difference in mind.

\subsection{Training and implementation}

All models use AdamW \citep{loshchilov2019decoupled} with cosine decay,
bfloat16 mixed precision, dropout 0.1, and batch size 32 except ISRUC (16) and
MentalArithmetic (64). The ImageNet/DINOv3 backbone is fine-tuned jointly with
a newly initialized linear head, and bias and normalization parameters receive
zero weight decay. Samples are converted to float32, clipped, and scaled
before folding (default $x\mapsto x/32$ after clipping to $[-1024,1024]$), and
the fold factor is fixed before training using Eq.~\ref{eq:p-geometry-rule}.
Dataset-specific learning rates, weight decay, epochs, EMA settings, and
augmentations are recorded in Appendix~\ref{app:recipes}; seeds, checkpoint
selection, and the final-test protocol follow the locked YAML recipes.

\paragraph{Ablation and control studies.}
We evaluate the interface through four complementary studies: \textbf{(1)
Cross-backbone consistency} tests ImageNet-pretrained EfficientNet-B0 and
DINOv3-pretrained ConvNeXt-Tiny and ViT-Small (Table~\ref{tab:unified-results};
Appendix~\ref{app:detailed-results}). \textbf{(2) Pretraining contribution}
uses matched pretrained-versus-random initialization under identical recipes
(Section~\ref{sec:pretraining-ablation}; Table~\ref{tab:vision-pretraining-ablation}).
\textbf{(3) Folding layout} compares phase-interleaved folding with contiguous
reshape; $P=1$ is the identity, so the informative comparison is $P>1$
(Section~\ref{sec:fold-reshape}; Appendix~\ref{app:fold-reshape}).
\textbf{(4) Fold-factor sensitivity} sweeps $P\in\{1,2,4,8\}$ on CHB-MIT,
TUEV, and MentalArithmetic (Appendix~\ref{app:p-ablation}); tuned factors use
validation data only.

\subsection{Evaluation protocol}

We compare \method{} with task-specific EEG models and EEG foundation models,
using published baseline results from \cbramod{} and \reve{} under the same
downstream data splits \citep{wang2024cbramod,elouahidi2025reve}. For \cbramod{}
on TUAB and TUEV, we report the variants whose pretraining excludes the target
corpus.

Binary tasks report balanced accuracy (BA), PR-AUC, and ROC-AUC; multiclass
tasks report BA, Cohen's $\kappa$, and weighted F1. Checkpoint selection uses
only validation PR-AUC or validation $\kappa$, while the remaining metrics are
test-time reporting metrics; published BAs are transcribed from the same
benchmark tables and are never mixed across runs. Each recipe is trained with
five seeds (42--46), and we report the arithmetic mean and population standard
deviation after selecting one checkpoint on validation and evaluating it once
on test; the executable protocol is included in the released code.
